\chapter{Appendix A: Participant Questionnaire and Responses}
\label{app:questionnaire}

This appendix presents the questionnaire that was created to evaluate game developer demographics and the 
approachability of academic notation, followed by the data generated from the aggregated responses and qualitative feedback

\section{Blank Questionnaire}
\label{sec:blank_questionnaire}

\subsubsection{Demographics \& Backgrounds}

\textbf{1. What is your primary background in game development?}

\begin{itemize}
    \item Indie
    \item Hobbyist
    \item Working in Industry
    \item Student
    \item I am in an adjacent field (Software Engineering, IT, etc)
\end{itemize}

\vspace{2em}

\textbf{2. What is your highest level of formal mathematics education?}

\begin{itemize}
    \item GCSE
    \item A-Levels
    \item Higher Education (Under/Postgraduate)
    \item Self Taught
    \item (Other)
\end{itemize}

\vspace{2em}

\textbf{3. How often do you refer to formal or informal academic media? (Research Papers, Conferences, Blog Posts, YouTube Videos etc.)}

\begin{itemize}
    \item Never
    \item Rarely
    \item Occasionally
    \item Frequently
\end{itemize}


\subsubsection{Paper Readers}

\textbf{4. [Notation] When reading a paper, how do you interact with with mathematical notation such as (e.g., $\sum$, $\prod$)}

\begin{itemize}
    \item I do not understand some of the notation used
    \item I rely on using AI to explain and give context to the notation used
    \item I skip past the maths and try look for code snippets
    \item I look up their definitions on external sources (e.g other papers, google, etc)
    \item I understand them comfortably
\end{itemize}

\vspace{2em}

\textbf{5. [Lexicon] How would you rate the average clarity of the language used in peer-reviewed papers vs. blogs or other sources}

\begin{table}[h]
\centering
\begin{tabular}{lcccc}
\toprule
Item & Dense & Confusing & Clear & Understandable \\
\midrule
Peer-Reviewed Papers & [ ] & [ ] & [ ] & [ ] \\
Industry blogs / other sources & [ ] & [ ] & [ ] & [ ] \\
\bottomrule
\end{tabular}
\end{table}

\vspace{2em}

\textbf{6. [Contextualisation] Which of the following elements add the most value when attempting to implement a complex algorithm from a paper (Select all that Apply)}

\begin{itemize}
    \item High level psuedo-code
    \item Fully implemented code snippets (C++, C\#, etc.)
    \item Step-by-step diagrams
    \item Formal writing
    \item All of the above
    \item None of the above
    \item (Other)
\end{itemize}

\vspace{2em}

\textbf{7. [Pedagogy] How often do you find that a paper attempts to TEACH you what a concept is, instead of just providing clarification}

\begin{itemize}
    \item Never
    \item Rarely
    \item Occasionally
    \item Frequently
\end{itemize}


\subsubsection{Barriers to Entry}

\textit{You answered that you rarely, if ever refer to academic literature, why?}

\vspace{1em}

\textbf{8. What are the primary reasons that you do not use academic literature for game development? (Select all that Apply)}

\begin{itemize}
    \item I feel like I do not need to use academic papers
    \item I do not know how to find relevant papers
    \item The papers and their contents are daunting (notation, context, etc.)
    \item Too much academic jargon
    \item I feel like there is a `disconnect' from the theory in the paper and the implementation in engine
    \item I prefer tutorials, blogs, or other sources
\end{itemize}

\vspace{2em}

\textbf{9. If academic literature consistently used functional code samples alongside their mathematical content, how likely would you be to use them?}

\begin{itemize}
    \item Unlikely
    \item Neither likely or unlikely
    \item Likely
\end{itemize}


\subsubsection{Code vs. Notation}

\textbf{10. Look at the following expression. It can be frequently found in literature about procedural generation. How confident are you that you could accurately rewrite this into a for loop in your preferred programming language WITHOUT looking it up?}

$$P(x) = \sum_{i=1}^{n} f(i)$$

\begin{itemize}
    \item No idea where to start
    \item Somewhat confident
    \item Confident
\end{itemize}

\textbf{11. When using native game engine features such as Unity's Perlin noise functions, how important is it to you that you understand the underlying concepts and mathematics behind them?}

\begin{itemize}
    \item Not important
    \item Somewhat unimportant
    \item Neutral
    \item Somewhat important
    \item Extremely important
\end{itemize}

\vspace{2em}

\textbf{12. Any further thoughts?}

\begin{center}
    (Text field)
\end{center}

\pagebreak

%%%%%%%%%%%%%%%%%%%%%%%%%%%%%%%%%%%%%%%%%%%%%%%%%%%%%%%%%%%%%%%%%%%%%%%%%%%%%%%%%%%%%%%%%%%%%%%%%%%%%%%%%%%%%%%%%%%%%%%%%%%%%

\section{Aggregated Responses and Qualitative Feedback}
\label{sec:questionnaire_responses}

This section presents the raw data counts and the qualitative feedback collected through the questionnaire detailed in Appendix \textit{A}.

\subsubsection{Demographics \& Backgrounds}

\textbf{1. What is your primary background in game development?}

\begin{itemize}
    \item Indie (2 responses)
    \item Hobbyist (1 response)
    \item Working in Industry (3 responses)
    \item Student (9 responses)
    \item I am in an adjacent field (Software Engineering, IT, etc) (0 responses)
\end{itemize}

\vspace{2em}

\textbf{2. What is your highest level of formal mathematics education?}

\begin{itemize}
    \item GCSE (8 responses)
    \item A-Levels (2 responses)
    \item Higher Education (Under/Postgraduate) (5 responses)
    \item Self Taught (0 responses)
    \item (Other) (0 responses)
\end{itemize}

\vspace{2em}

\textbf{3. How often do you refer to formal or informal academic media? (Research Papers, Conferences, Blog Posts, YouTube Videos etc.)}

\begin{itemize}
    \item Never (0 responses)
    \item Rarely (1 response)
    \item Occasionally (4 responses)
    \item Frequently (10 responses)
\end{itemize}


\subsubsection{Paper Readers}

\textbf{4. [Notation] When reading a paper, how do you interact with with mathematical notation such as (e.g., $\sum$, $\prod$)}

\begin{itemize}
    \item I do not understand some of the notation used (1 response)
    \item I rely on using AI to explain and give context to the notation used (1 response)
    \item I skip past the maths and try look for code snippets (2 responses)
    \item I look up their definitions on external sources (e.g other papers, google, etc) (7 responses)
    \item I understand them comfortably (4 responses)
\end{itemize}

\vspace{2em}

\textbf{5. [Lexicon] How would you rate the average clarity of the language used in peer-reviewed papers vs. blogs or other sources}

\begin{table}[h]
\centering
\begin{tabular}{lcccc}
\toprule
Item & Dense & Confusing & Clear & Understandable \\
\midrule
Peer-Reviewed Papers & 4 & 8 & 1 & 1 \\
Industry blogs / other sources & 1 & 0 & 9 & 4 \\
\bottomrule
\end{tabular}
\end{table}

\vspace{2em}

\textbf{6. [Contextualisation] Which of the following elements add the most value when attempting to implement a complex algorithm from a paper (Select all that Apply)}

\begin{itemize}
    \item High level psuedo-code (5 responses)
    \item Fully implemented code snippets (C++, C\#, etc.) (5 responses)
    \item Step-by-step diagrams (4 responses)
    \item Formal writing (3 responses)
    \item All of the above (6 responses)
    \item None of the above (0 responses)
\end{itemize}

\vspace{2em}

\textbf{7. [Pedagogy] How often do you find that a paper attempts to TEACH you what a concept is, instead of just providing clarification}

\begin{itemize}
    \item Never (1 response)
    \item Rarely (9 responses)
    \item Occasionally (4 responses)
    \item Frequently (0 responses)
\end{itemize}


\subsubsection{Barriers to Entry}

\textit{You answered that you rarely, if ever refer to academic literature, why?}

\vspace{1em}

\textbf{8. What are the primary reasons that you do not use academic literature for game development? (Select all that Apply)}

\begin{itemize}
    \item I feel like I do not need to use academic papers (1 response)
    \item I do not know how to find relevant papers (0 responses)
    \item The papers and their contents are daunting (notation, context, etc.) (1 response)
    \item Too much academic jargon (1 response)
    \item I feel like there is a `disconnect' from the theory in the paper and the implementation in engine (1 response)
    \item I prefer tutorials, blogs, or other sources (1 response)
\end{itemize}

\vspace{2em}

\textbf{9. If academic literature consistently used functional code samples alongside their mathematical content, how likely would you be to use them?}

\begin{itemize}
    \item Unlikely (0 responses)
    \item Neither likely or unlikely (1 response)
    \item Likely (0 responses)
\end{itemize}


\subsubsection{Code vs. Notation}

\textbf{10. Look at the following expression. It can be frequently found in literature about procedural generation. How confident are you that you could accurately rewrite this into a for loop in your preferred programming language WITHOUT looking it up?}

$$P(x) = \sum_{i=1}^{n} f(i)$$

\begin{itemize}
    \item No idea where to start (8 responses)
    \item Somewhat confident (3 responses)
    \item Confident (4 responses)
\end{itemize}

\vspace{2em}

\textbf{11. When using native game engine features such as Unity's Perlin noise functions, how important is it to you that you understand the underlying concepts and mathematics behind them?}

\begin{itemize}
    \item Not important (0 responses)
    \item Somewhat unimportant (3 responses)
    \item Neutral (6 responses)
    \item Somewhat important (4 responses)
    \item Extremely important (2 responses)
\end{itemize}

\vspace{2em}

\textbf{12. Any further thoughts?}

\textit{Note: Due to these being direct quotations taken from the questionnaire, these answers are presented verbatim. They have not been modified for spelling, grammar, or other reasons}

\begin{itemize}
    \item \textit{Respondent 6:} ``i find myself wanting to understand underlying concepts mostly for optimisation reasons, if i understand the concepts i can sometimes use a cheaper approximation to a given algorithm''
    \item \textit{Respondent 12:} ``This questionaire would be well served by asking for an equation to be solved, then asking the exact steps taken to solve it as you simply remove five questions and aee given a full mental landscape runthrough in one text box. I suppose especially regarding the utilisation of features like “Perlin noise” it’s worth asking firstly what the game designeer or person believes perlin noise to be and how they cane to that knowledge (for example, as an artist I use chat GPT advanced to transplate english programatical requirwmwnts to C++ with around a 90\% efficiency rate which is crucial knowledge for a dissertation)''
\end{itemize}

